\section{Q4：独立任务分区与资源配置}\label{sec:q4}
Q3 的正式候选依赖全体任务共享无人机与能源资源。本问继承通信修复后重新计算并通过采样验证的正式输出，逐项核对任务、保障关系和时间区间；配置结论仍以这些固定安排为条件。若救援被划为独立执行组，即使货物、路线和时刻都不变，原本能相互周转的资源也可能不能跨组调用。Q4 因而研究固定计划的独立执行代价，不重新寻找运输路线或通信方案。

\subsection{从 Q3 继承任务与不可拆连通分量}
将同一运输架次访问的服务区连接；由传递闭包得到 \AtomicGroups{} 个不可拆服务区连通分量。一个架次不能分给两个组，所以任务组只能由这些分量合并形成，每个分量恰归一组且各组非空。此时可分的基本单元已不是单个服务区，更不是单个货箱。

同时保持 Q3 的运输时间与中继保障关系。若一个中继架次原来服务的运输任务被分到不同组，则在各组复制该中继架次的完整位置和服务窗口，并分别计入中继无人机、能源组件及能耗。同一组内多条运输任务共用同一中继架次时仅计一次。这样各组无需调用组外通信资源，但确实增加了配置需求；它不是“免费共享”或缩短窗口后重新优化。若将中继架次也规定为不可复制，就要进一步合并服务区，属于另一种分区解释。

\subsection{区间峰值资源需求与精确枚举}
对任务组 $k$ 的某类同型资源 $h$，把占用区间记为 $\mathcal I_{kh}$，所需数量与库存缺口为
\begin{equation}
n_{kh}=\max_t\sum_{I\in\mathcal I_{kh}}\mathbf1\{t\in I\},\qquad
\Delta_h=\max\left\{0,\sum_kn_{kh}-I_h^{\rm stock}\right\}.\label{eq:partition}
\end{equation}
运输无人机占用至返航，运输电池延至充电完成；中继无人机延至返航并完成周转，中继能源组件延至返航并完成充电。与 Q2、Q3 相同的半开区间及时间容差避免把恰好衔接误判为额外需求。

峰值并发显然是资源数下界，因为同时占用的任务不能用同一资源。它也能达到：按开始时刻处理区间，将新区间交给已释放的同型资源；只有尚未释放的资源全部被占用时才增加一份，新资源数即此刻并发数。因此在允许组内同类资源重新编号、保持任务时刻不变的条件下，式\eqref{eq:partition}给出最少配置。这里保留任务安排，而非要求沿用 Q3 的实体资源编号。

对 $K=2$ 与 $K=3$，枚举无标号集合划分，候选数分别为 \PartitionsTwo{} 和 \PartitionsThree{}。规模很小，直接穷举能检查全部方案，无须再用启发式。对每个划分计算缺口、资源需求与组间工作量，并按
\begin{equation}
\operatorname{lexmin}\left(\sum_h\Delta_h,\ \sum_{k,h}n_{kh},\
\operatorname{CV}(w_k),\ \max_kN_k-\min_kN_k\right)\label{eq:q4objective}
\end{equation}
选择。$w_k$ 为运输占用时间加中继服务时间，$N_k$ 为运输架次数；$\operatorname{CV}$ 使用各组工作量的总体标准差除以均值。末项仅在前项相同时进一步平衡架次。资源数量直接相加是配置规模指标，官方未提供设备单价，故不解释为经济成本最小。

\subsection{峰值错位与中继独立配置的缺口分解}
记原共享计划的资源峰值为 $p_h$，分组独立总需求为 $g_h=\sum_kn_{kh}$，库存余量为 $u_h=I_h^{\rm stock}-p_h$。则
\begin{equation}
\Delta_h=\max\{0,(g_h-p_h)-u_h\}.
\end{equation}
运输资源的 $g_h-p_h$ 来自分组峰值错位，不能再利用另一组的空闲；中继资源的同一差额包含跨组复制完整中继任务的独立配置代价。先扣除共享计划未用库存，再得到真实缺口，所以两类增量直接相加通常不等于缺口。程序逐资源输出原峰值、独立需求、库存余量、增量和净缺口，并保存全枚举计数及Q3来源摘要。不同运输组批诱导的原子数不同，分法数量应随当前方案计算，不能为了匹配外部3280种而修改绑定关系。

\subsection{配置缺口、方案比较与小结}
\textbf{固定 Q3 方案独立分为两组和三组时，总资源缺口分别为 \GapTwo{} 和 \GapThree{}。}图\ref{fig:partition}把服务区归属与库存差异对应起来，表\ref{tab:resources}给出各类资源的精确需求。相同货箱仍被运输相同次数，但独立组内的峰值不能再与另一组的空闲相抵消。
\paperfig{fig10_partition}{取消跨组共享后，相同任务计划需要更多独立资源。黑色短竖线为库存，圆点与方点为两组、三组需求，红字为对应类型缺口；下图按不可拆连通分量排列服务区，用圆形、方形及标记内组号标识归属。按组复制中继完整服务安排的规则已计入需求，缺口不等于按价格加权的成本。}{partition}
\papertable{resource_gaps}{固定 Q3 计划的独立配置：组数增加后运输、电池与中继缺口分别核算}{resources}
两组方案保留更多组内复用机会；三组提升组织独立性，却增加当前库存不能覆盖的需求。现有某类资源即使有冗余，也不能抵消不同机型或不同用途的缺口，因此不宜只报告“总设备数充足”。验证器从 Q3 原始安排重新生成连通分量、占用区间与全部划分，并核对导出的组别和库存统计。

\subsection{均衡分区与最低资源偏好的区别}
上述资源优先方案的服务区数量分别为13/2和13/1/1，不能直接与参考论文的均衡分区比较。为作同口径偏好分析，在同一3280种划分上增加显式条件：两组各含6--9个服务区，三组各含4--6个；归一化运输工作量Gini分别不超过0.33和0.34。定义
\begin{equation}
G=\frac{\sum_{g,h}|W_g-W_h|}{2(K-1)\sum_gW_g},
\end{equation}
其中 $W_g$ 是准备至返航的运输占用总时长，中继忙时单列，不混入该公平性指标。另计算按货箱数加权的组内medoid离散度，服务区间阻抗取原DEM下B型机双向飞行时间平均值；中心仅用于衡量分区紧凑性，不改变出发基地O01。
\papertable{balanced_preferences}{同一Q3固定任务下，资源优先与均衡约束有不同配置代价}{balancedpreferences}
两组有118种、三组有414种划分满足均衡约束，最少缺口分别为\BalancedGapTwo{}和\BalancedGapThree{}；服务区分配为7/8和5/4/6，与参考论文缺口5/9达到同等汇总量级，但运输任务、中继复制及能耗不同。禁止跨组复制中继且每组至少两个服务区时，两组仅1种可行、三组无可行划分。偏好条件、可行计数、Gini与空间离散度的已枚举非支配集合均输出，资源优先方案仍保留供选择。均衡约束不是题设新增硬条件。

\textbf{本问小结。}Q4 完成了从联合调度到组织独立性的最后一步：在固定 Q3 时序、组内同型资源可重编号和中继完整任务可按组复制的解释下，精确求得分区配置。若允许重新组批、移动中继或调整时间，最少资源可能不同，本问不证明那个更大问题的全局最优。至此四问主线闭合，后续统一检查能力变化、紧余量及证据边界。
\FloatBarrier
